% from ega4-12-fr.tex
\oldpage[IV-3]{179}
  mais en vertu de (ii) et des remarques précédentes, cet ensemble est
  intersection de deux ouverts de $X$, donc est ouvert dans $X$. C.Q.F.D.
\end{enumerate}
\end{proof}

\begin{corollary}[12.1.7]
\label{IV.12.1.7-fr}
Soit $f:X\to Y$ un morphisme localement de présentation finie. Alors
l'ensemble des points $x\in X$ où $f$ possède l'une des propriétés suivantes
(6.8.1)~:
\begin{enumerate}[label=(\roman*)]
  \item vérifier la propriété $(S_k)$~;
  \item être de coprofondeur $\leq n$~;
  \item être de Cohen-Macaulay~;
  \item être régulier (ou lisse, ce qui revient au même)~;
  \item être normal~;
  \item être réduit~;
\end{enumerate}
est ouvert dans $X$.
\end{corollary}

\begin{proof}
En effet, il résulte de (11.3.1) que l'ensemble des $x\in X$ où $f$ est plat
est ouvert. On peut donc se borner au cas où $f$ est plat, et alors le
corollaire résulte de (12.1.1, (iv), (vi) et (vii)) et de (12.1.6, (i), (ii)
et (iv)).
\end{proof}

\begin{remarks}[12.1.8]
\label{IV.12.1.8-fr}
\begin{enumerate}[label=(\roman*)]
  \item Dans (12.1.6, (ii)), on ne peut supprimer l'hypothèse que $x$
  n'appartienne à aucun cycle premier associé immergé de $X_{f(x)}$. C'est ce
  que montre l'exemple (5.12.3), où l'on prend $X=\operatorname{Spec}(A)$, $Y$
  étant le spectre de l'anneau local $A_0$ de $k[T]$ correspondant à l'idéal
  premier $(T)$ et le morphisme $X\to Y$ correspondant à l'homomorphisme
  injectif $A_0\to A$, déduit par localisation et passage au quotient de
  l'injection $k[T]\to k[T,U,V]$~; ce morphisme est plat puisque $A$ est un
  $A_0$-module sans torsion ($\mathbf 0_{\mathrm I}$, 6.3.4). Ici la fibre
  $X_y$ au point fermé $y$ de $Y$, égale à $\operatorname{Spec}(A/tA)$, est
  irréductible, de dimension 1 et vérifie la condition $(R_0)$ géométrique
  puisque l'anneau local en son point générique est un corps. Par contre, au
  point générique $y'$ de $Y$, la fibre $X_{y'}$ a deux composantes
  irréductibles de dimensions respectives 0 et 1, et celle de dimension 0
  n'est pas réduite, donc $X_{y'}$ ne vérifie pas la condition $(R_0)$.
  \item Dans (12.1.6, (ii)), on ne peut pas non plus supprimer l'hypothèse que
  $X_{f(x)}$ est équidimensionnel au point $x$. On le voit ici sur l'exemple
  (5.12.6) avec $X=\operatorname{Spec}(A)$, $Y=\operatorname{Spec}(A_0)$, où
  $A_0$ est défini comme dans (i), le morphisme $X\to Y$ provenant encore de
  l'injection $k[T]\to k[T,U,V,W]$ par localisation et passage au quotient, et
  étant plat puisque $A$ est un $A_0$-module sans torsion
  ($\mathbf 0_{\mathrm I}$, 6.3.4). La fibre $X_y$ au point fermé $y$ de $Y$
  est réduite (donc vérifie $(S_1)$) et vérifie $(R_1)$, mais a deux composantes
  irréductibles de dimensions 2 et 1, tandis que la fibre $X_{y'}$ au point
  générique $y'$ de $Y$ ne vérifie pas $(R_1)$.
\end{enumerate}
\end{remarks}
\subsection{Propriétés locales et globales des fibres d'un morphisme propre,
plat et de présentation finie}
\label{subsection:IV.12.2-fr}

\begin{theorem}[12.2.1]
\label{IV.12.2.1-fr}
Soient $f:X\to Y$ un morphisme \emph{propre} de présentation finie,
$\mathcal F$ un $\mathcal O_X$-Module $f$-plat et de présentation finie,
$\Phi$ une partie finie de $\mathbf Z\cup\{\pm\infty\}$, $k$ un entier. Les
parties suivantes de $Y$ sont ouvertes~:
\begin{enumerate}[label=(\roman*)]
  \item L'ensemble des $y\in Y$ tels que l'ensemble des dimensions des cycles
  premiers associés à $\mathcal F_y$ soit contenu dans $\Phi$.
  \item L'ensemble des $y\in Y$ tels que l'ensemble des dimensions des
  composantes irréductibles de $\operatorname{Supp}(\mathcal F_y)$ contienne
  $\Phi$.
  \item L'ensemble des $y\in Y$ tels que tous les cycles premiers associés à
  $\mathcal F_y$ aient une même dimension égale à $k$.
  \item L'ensemble des $y\in Y$ tels que $\mathcal F_y$ n'ait pas de cycle
  premier associé immergé et que l'ensemble des dimensions des composantes
  irréductibles de $\operatorname{Supp}(\mathcal F_y)$ soit égal à $\Phi$.
  